\section{Evaluation methodology}
\label{ch:methodology}

\subsection{Purpose and comparison design}

The controlled experiment asks whether two evidence-selection choices change responses to institutional questions. Production retrieval and prompt construction are reused, while generation runs through a controlled hosted client, so the conclusions apply to this narrower execution path.

The main sample contains 192 questions, divided equally between email and production frames, preceded by a separate 20-question email pilot. Sampling uses seed 24905307 and retained manifests. The email questions had already been examined in the component study, and the production sample consists of selectively recorded latest turns. Chapter~\ref{ch:data} explains why neither frame is a complete population of student needs.

The policies cross two factors: the document score floor and whether an all-below-floor fallback is enabled. The FAQ floor stays at 0.10. The document floor is 0.10 in the shared-threshold policies and 0.00 in the separate-threshold policies. G0 and S0 disable fallback; G1 and S1 enable it. 

\begin{table}[htbp]
\centering
\caption{Frozen retrieval policies. The source-display threshold remains 0.30 and the final candidate budget remains ten in every policy.}
\label{tab:thesis-policies}
\begin{tabular}{lrrc}
\toprule
Policy & FAQ floor & Document floor & Fallback \\
\midrule
G0 & 0.10 & 0.10 & Off \\
G1 & 0.10 & 0.10 & On \\
S0 & 0.10 & 0.00 & Off \\
S1 & 0.10 & 0.00 & On \\
\bottomrule
\end{tabular}
\end{table}

The primary contrast is S0 minus G0, which changes the document floor with fallback disabled. S1 minus G1 examines the same change with fallback enabled. G1 minus G0 and S1 minus S0 examine fallback under each threshold, and the interaction asks whether the fallback effect differs between them.

The fallback retains the single highest-scoring candidate only when no candidate from either corpus survives the floor; it is not applied per corpus. If a document already survives under S0, enabling fallback may change nothing. A separate application recovery path for reranker failure, which uses dense ordering, is not the policy treatment: an occurrence is counted as an execution error.

\subsection{Production retrieval capture and replay}

The inspected API source files match local files by SHA-256. The evaluation imports the production retrieval function and captures candidate order, source content, dense and lexical retrieval results, and neighboring document chunks. Each policy thus sees the same captured evidence.

Database capture uses server-enforced read-only repeatable-read transactions. Each gives a consistent view within its batch, but separate batches may observe different moments, so the study checks repeated source content across batches instead of claiming one snapshot. The retained candidates contain 9,329 occurrences covering 681 unique identifiers, and no repeated identifier has conflicting text for reranking. This supports consistency of the captured material, not of uncaptured records or future live queries.

BGE-M3 query embeddings and BGE-reranker-v2-m3 scores are computed through isolated imports of the deployed GPU modules~\cite{chen2024m3,baaiBgeRerankerV2M3} and saved for reuse. Link reranking is captured separately, so the threshold experiment cannot change link selection.

Replay applies two controlled changes to the production function's abstract syntax tree: the score-floor comparison and the fallback condition. Candidate budgets, ordering, FAQ priority, neighboring-chunk expansion and source-display behavior remain fixed. Raw retrieval allows up to 31 dense candidates per corpus before merging and any permitted lexical expansion; the final budget is ten. Scores equal to the floor remain eligible.

The context floor for documents can be 0.10 or 0.00, but the source-display floor remains 0.30. A passage may therefore enter the generation context without becoming a visible source card; more model context need not mean more evidence shown to the student.

\subsection{The evaluated generation boundary}

The primary query mode is raw, without paid query rewriting or hypothetical-document generation. The experiment also excludes conversation history, remote MCP tools, individual committee transactions, HTTP authorization, quota accounting, reverse-proxy behavior and browser interaction. The run is therefore not a test of the complete deployed agent.

The application's \texttt{handle\_chat} function constructs the system and user prompts, retaining interface instructions, empty-context wording, useful links and a fixed date of 07.09.2026. The evaluator intercepts the call before production provider dispatch and sends the captured prompt through a bounded Chat Completions client, whereas production uses a Responses/agent path; the two transports are not shown to have identical defaults or tool behavior.

Luna is the requested model alias \texttt{gpt-5.6-luna}, accessed through an OpenAI-compatible gateway that reports \texttt{gpt-5.6-luna-}\allowbreak\texttt{2026-07-09}; this identity is recorded but not independently certified. Generation uses reasoning effort none, a 1,024-token output limit, no automatic retry and the gateway's default sampling settings. The application's temperature and top-$p$ values are retained as metadata but not sent.

\subsection{Prompt identity, reuse and execution accounting}

When question, prompt bytes and generation settings are identical across policies, the evaluator reuses one generated answer. Unique requests are shuffled with a fixed seed. Reuse prevents artificial policy differences caused only by independent samples.

The four policies produce 768 nominal main question--policy records. Sixteen input rejections leave 752 eligible records, which map to 354 unique main generation calls (195 email, 159 production); reuse accounts for the remaining 398. The pilot adds 39 unique calls, kept separate from the main estimates.

Execution outcomes are preserved: a rejected input, incomplete output, transport problem or invalid assessment cannot become a success by leaving the denominator. Reused answers remain linked to every policy they represent, so one assessment problem can appear in several policy records; the results distinguish unique failed judgments from the nominal outcomes they map to.

Capture and replay block the application's HTTP dispatch and do not run migrations, save conversations, change quotas or initialize production analytics, so no synthetic questions enter the historical usage observations reported later.

\subsection{Reference construction before answer assessment}

Before candidate answers are evaluated, Luna builds one reference per question, shared by all policies. The reference prompt sees the question and a bounded evidence pool but not the answers to be judged. Because the same model family builds references, generates answers and judges them, model-specific errors are not eliminated.

The source pool alternates between up to eight FAQ records and eight document chunks, ordered by score within each type. Each source contributes at most its first 1,400 characters, within a total of 12,000. The pool can omit a necessary provision or truncate a condition, and it favors evidence already found by retrieval, which limits the reference's independence from the system being studied.

The reference records answerability, required facts with source identifiers and quotations, required action, currentness and uncertainty. Code checks whether each quotation occurs in its named source; facts failing the check remain in the raw output but are removed from the checked list, and the reference is marked uncertain. This detects fabricated or misattributed quotations but does not test whether a quotation entails the fact. When references and judgments are combined, an otherwise successful supported-answer or safe-routing label becomes uncertain if the reference fails validation; negative and incomplete judgments remain as recorded.

References with no retained facts can pass these checks. For a greeting or clarification an empty target may be reasonable, but for a substantive question it leaves factual completeness unchecked. The results therefore report empty references explicitly, with a separate post-hoc sensitivity analysis that does not replace the frozen primary measure.

\subsection{Answer assessment and outcome categories}

The primary judge receives one anonymous candidate answer, its question, the common reference and the exact captured generation prompt, but not the policy label or earlier judgments. The executed output schema is compact: \texttt{overall}, \texttt{unsupported\_material\_claim}, \texttt{all\_required\_facts\_covered}, \texttt{citation\_supported}, \texttt{rationale} and \texttt{evidence\_quotes}. The initial design proposed per-claim arrays, which were not implemented (Section~\ref{sec:method-deviations}).

The overall category distinguishes supported answering, useful safe routing, unsupported material claims, incompleteness, unnecessary abstention, uncertainty and failure. A material claim concerns consequential information such as eligibility, fees, deadlines, procedure or the claimed completion of an institutional action. A true unsupported-material-claim flag requires the unsupported category. The flag can also be unspecified, which affects the denominator for flag-level agreement. Structural validation checks required fields, permitted categories and internal consistency, rejecting, for example, a successful label combined with an unsupported material claim. Passing it means the record is interpretable, not that its judgment is correct.

Supported answering requires the reference's facts and conditions. Useful routing includes necessary clarification, advice to verify unresolved information, or referral to the appropriate service without inventing a completed private transaction. The separate categories remain reported because they show what the student actually receives.

For example, an AI-authored control asks what to submit for HPC credit recognition, where, and the paper application fee according to a frozen FAQ. The candidate gives only the 50-denar fee, omitting the printed certificates, the recognition application and submission to the archive or Student Services. The fee is supported but the answer is incomplete. The control illustrates the rubric and is not a user response or current fee advice.

\subsection{Calibration and assessment sensitivity}

The evaluation includes 18 source-based AI-authored controls covering contradictions, omissions, incorrect citations, fabricated transactions, justified refusal and applicability limits. A protocol addendum, recorded before answer generation, adds twelve Macedonian renditions arranged as six Cyrillic/Latin pairs. These are development checks, not validated gold-standard examples or a study of student language preferences.

The 20-case pilot checks that references, generation, judging and analysis execute under the planned settings; it does not contribute to the main estimates and is evidence about the pipeline, not judge accuracy. No threshold is selected by searching pilot or main answer outcomes.

Assessment is repeated with identical prompts for a fixed subset of 48 unique answers. The same prompts and references are also sent to a pinned local Qwen3-8B Q4\_K\_M model~\cite{yang2025qwen3}, with temperature zero, a 1,024-token limit, a 90-second timeout and no retry. Repeated Luna judgments examine within-assessor consistency; Qwen judgments examine variation across assessors under matched evidence. Neither is human validation.

After substantial cross-model disagreement was observed, Qwen was also run on the unchanged thirty control prompts to test whether the two models read the rubric similarly. This post-hoc diagnostic does not change the primary labels or imply that the harsher model is correct.

\subsection{Paired outcome estimates and uncertainty}

Let $Y_{i,p}$ equal one when question $i$ under policy $p$ retains either a supported-answer or useful-routing label, and zero otherwise. For each frame, the main estimate is
\begin{equation}
\widehat{\Delta}_{\mathrm{floor}}=\frac{1}{96}\sum_{i=1}^{96}\left(Y_{i,S0}-Y_{i,G0}\right).
\label{eq:thesis-primary}
\end{equation}
The reported percentage-point difference is $100\widehat{\Delta}_{\mathrm{floor}}$. Input rejections, invalid judgments and uncertainty contribute zero to this endpoint while remaining visible in their own categories. A conditional analysis of only assessable records would answer a different question and needs a separately stated denominator.

The interaction is the average of $(Y_{i,S1}-Y_{i,S0})-(Y_{i,G1}-Y_{i,G0})$. Question groups are resampled 20,000 times with seed 24905310, keeping each question's four policy outcomes together, and percentile intervals are reported~\cite{efron1993bootstrap}. This preserves the paired structure and the dependence created by reuse. Email and production estimates are reported separately; any combined summary describes the balanced sample, not deployed usage.

The intervals quantify variation across sampled questions conditional on the saved AI labels. They do not correct a consistently wrong judge, a missing reference condition or a biased question frame, and an interval spanning zero does not establish equivalence.

Agreement is reported as counts, percentages and Cohen's kappa~\cite{cohen1960kappa}, which compares observed agreement with that expected from the assessors' marginal frequencies. High raw agreement can coexist with lower kappa when one label is common~\cite{feinstein1990kappa}. Missing or unspecified flags require separate paired denominators. Reporting category and material-claim agreement together shows whether assessors disagree about the overall response or about a consequential factual issue.

\subsection{Retrieval diagnostics, timing and cost}

To interpret policy effects, the analysis reports empty contexts, source mix, context length and identical context/source identifiers across policies. Empty-to-nonempty changes indicate observed fallback exposure, not an instrumented count of every fallback branch. Occurrence of reference quotations in context measures availability of the retained quote strings, not semantic recall.

Latency and token usage are recorded for unique calls; replay timing, hosted generation time and historical deployment duration cover different intervals and are not interchangeable. Historical events from PostHog, the deployment's product-analytics service, are a separate dated observational source, not outcome labels.

The inference client reserves costs at conservative planning rates of USD~2 per million input tokens and USD~12 per million output tokens. Its experiment ledger has an USD~80 reservation ceiling and a 2,500-request cap; a separate identity diagnostic has its own small reservation. Failed requests retain reservations. Actual gateway billing and total infrastructure cost were not available.

\subsection{Prompt comparison}
\label{sec:method-prompt}

A second experiment compares the production prompt with a revised prompt on an expanded collection: a read-only snapshot of 74 FAQs, sixty documents and 2,391 chunks. FAQ text matches the retrieval-policy experiment, but the document set is larger.

A 32-question development sample, balanced between email and production sources and spread across the reference states of the retrieval-policy experiment, compares current retrieval, AI-selected reference evidence and no evidence. It motivated a prompt-only revision, frozen before the comparison: check dates and applicability instead of automatically favoring FAQs or higher-ranked text, preserve needed procedural steps, and state material gaps. The revision adds no model calls and leaves retrieval unchanged.

The comparison uses the first forty emails of a shuffled 64-question reserve, excluding exact and near duplicates of evaluated questions; semantic independence from those questions is not proved. Two AI passes prepare and review references before generation, selecting source lines to avoid transcription errors. All forty references pass exact-text checks (34 partial, three insufficient, three nonfactual); they are not institution-approved answers.

Both conditions use production retrieval with fallback, useful links, Luna, gateway sampling defaults and a 1,536-token cap. Strict completion requires every reference fact and behavior, no unsupported claim or citation error, and no reference concern; empty factual or mixed references cannot succeed. This endpoint is stricter than, and not comparable with, the retrieval-policy assistance measure. Paired differences use 20,000 bootstrap resamples, complemented by an exact paired test on the discordant cases.

\subsection{Protocol and deviations}
\label{sec:method-deviations}

The retrieval-policy experiment was frozen on 7 September 2026, before any evaluated answer was generated. The frozen protocol fixed the sample and its seed, the four policies, the primary contrast S0 minus G0, the reference and judgment procedures, the outcome categories and the analysis. Sample, reference and generation manifests record content hashes of the frozen inputs; the reference freeze alone covers 220 files. A protocol addendum, also recorded before answer generation, added the twelve Macedonian script controls and the input-rejection outcome for the four overlong questions; it changed no threshold or rubric. The prompt comparison had its own protocol, written before its answers were generated, and the revised prompt was frozen before the comparison ran.

The following analyses and changes were made after outcomes had been observed. They are reported as such, and none replaces a primary result:
\begin{itemize}
  \item the wording of the research questions was refined to clarify the effectiveness focus, without changing the frozen hypotheses, policies, outcomes or labels;
  \item the empty-reference sensitivity analysis, prompted by the number of references without retained facts;
  \item the Qwen run on the thirty control prompts, prompted by the observed cross-model disagreement;
  \item the offline diagnostic of rejected reference quotations;
  \item in the prompt comparison, the reassessment of all answer pairs with the full generation evidence, revised scope guidance and a 1,024-token judgment cap, prompted by a deterministic check showing that the initial judge lacked part of the generator's evidence.
\end{itemize}
The executed assessment also differs from the initial design: it uses the compact judgment schema rather than per-claim arrays, repeats 48 rather than 24 assessments, and omits a proposed fifth, tuned-threshold policy.

The retained records (sampling manifest, captured retrieval candidates, exact prompts, raw answers, all judgments and the analysis code) allow every reported count and interval to be recomputed offline without new inference. The private questions and answers are not released, and calling the hosted model again would not necessarily reproduce the same answer text. The retrieval-policy experiment lacks a complete database snapshot; the prompt comparison retains one.
